\section{Introduction}

Coverage-guided fuzzing has become the de facto approach for discovering vulnerabilities in operating system kernels. However, existing kernel fuzzers often encounter a fundamental bottleneck: after prolonged execution, the rate of coverage growth slows dramatically or even stagnates, leading to a plateau in vulnerability discovery. This phenomenon is particularly evident in long-running fuzzing campaigns such as Google Syzbot, where coverage gains in the later stages (after one month or more) become increasingly marginal. The root cause lies in the inherent limitation of coverage feedback alone, which struggles to guide the fuzzer toward deep program states where complex vulnerabilities reside. For large and intricate codebases like the Linux kernel, coverage feedback fails to capture the subtle behavioral differences triggered by distinct inputs even when traversing the same execution paths.

The significance of this problem is well recognized in recent research. First, mainstream kernel fuzzers such as Syzkaller exhibit diminishing returns in coverage growth during long-term campaigns. Empirical data from Google Syzbot and our own experiments confirm that coverage increases slow significantly after several weeks of continuous fuzzing. Second, relying solely on coverage feedback makes it difficult to efficiently discover deep vulnerabilities. Prior work has shown that coverage alone cannot effectively guide exploration of complex code paths, causing many bugs to remain undetected. Our experiments further reveal that discovering certain deep vulnerabilities using coverage-only approaches may require hundreds of days of execution. Third, crash samples accumulated over time contain rich vulnerability patterns that can provide valuable guidance for fuzzing. Research in hardware fuzzing domains has demonstrated the potential of leveraging historical crash data to improve fuzzing efficiency. Similarly, long-running campaigns such as Syzbot have amassed thousands of crash samples, and studies have identified recurring patterns among these crashes. Moreover, our experiments show that even crash samples from older kernel versions can assist in discovering vulnerabilities in newer versions, highlighting the cross-version relevance of such patterns.

Several existing approaches attempt to address this challenge, yet each has notable limitations. One category focuses on optimizing mutation strategies to accelerate coverage exploration. For example, recent techniques such as Mock and Snowplow refine mutation algorithms to reach certain code regions more quickly. While these methods achieve short-term improvements, they remain fundamentally constrained by coverage feedback and ultimately suffer the same stagnation when coverage growth plateaus over extended periods. Another category leverages external information, such as static analysis or large language models, to generate initial templates that improve code path reachability. Works like SyzDescribe and KernelGPT demonstrate the utility of synthesizing test templates from external sources. However, these approaches primarily enhance the initial phase of fuzzing and provide little benefit during the sustained, long-term execution phase. A third category attempts to incorporate manually crafted vulnerability patterns into the fuzzing process. For instance, ACTOR employs handcrafted patterns to guide fuzzing toward specific bug types. Although this approach moves beyond pure coverage feedback, manual pattern construction is labor-intensive and inherently limited to covering well-known, common vulnerability classes, leaving deeper and less obvious bugs unaddressed.

The fundamental technical deficiency shared by these methods is their exclusive reliance on coverage feedback, coupled with insufficient exploitation of historical crash data. Coverage feedback becomes less informative as the fuzzer saturates easily reachable code regions, resulting in diminished guidance for input generation. Meanwhile, the accumulated crash samples contain valuable information about error traces and execution paths that remain largely untapped. For complex targets like operating system kernels, different inputs can induce distinct behaviors even along identical code paths, and crash samples capture subtle interaction patterns that coverage metrics miss entirely.

Our key insight is to leverage historical crash samples collected by tools such as Syzbot to extract and generalize vulnerability patterns, which can then guide the fuzzer to generate more effective test inputs. By combining crash-derived patterns with coverage feedback, our approach maintains effective guidance even when coverage growth stagnates, thereby enabling continuous vulnerability discovery and accelerating the detection of deep bugs. This insight fundamentally differs from existing work in that we move beyond pure coverage feedback by integrating crash sample analysis. While prior methods depend solely on coverage metrics, our approach remains effective during prolonged campaigns when coverage plateaus. The feasibility of this insight is supported by several factors. First, large-scale fuzzing campaigns such as Google Syzbot have accumulated extensive crash datasets that serve as a foundation for pattern extraction. Second, emerging techniques in crash analysis and recent work in hardware fuzzing have demonstrated the viability of extracting vulnerability patterns from crash data. Furthermore, large language models can assist in abstracting and generalizing concrete parameter values found in crash samples into reusable patterns. We represent these generalized patterns using a Trie-based data structure, enabling random walks over the structure to generate novel test inputs. Finally, our method does not abandon coverage feedback entirely but rather combines it with crash-derived guidance to leverage the strengths of both approaches at different stages of the fuzzing campaign. Experimental results confirm that our method significantly improves vulnerability discovery efficiency, particularly after coverage growth has slowed.

Realizing this insight, however, presents two major technical challenges. The first challenge concerns the effective extraction of vulnerability patterns from crash samples. Syzbot and similar tools produce thousands of crash samples containing rich structural information along with substantial noise from specific parameter values. Additionally, the patterns within these samples exhibit a long-tail distribution, complicating extraction efforts. Existing solutions fall short in this regard. One approach, exemplified by ReFuzz in the hardware fuzzing domain, directly reuses previous test cases for mutation. However, this strategy is ineffective for kernel fuzzing due to significant code changes across kernel versions, rendering old test cases obsolete. Another approach, represented by tools such as MoonShine and KSG, employs program analysis to extract templates from seeds based on coverage feedback. While useful for expanding code coverage, these methods focus on coverage rather than the system call structure and parameter relationships critical to vulnerability patterns.

The second challenge involves effectively integrating the extracted patterns into the fuzzing mutation process. The integration mechanism must enable the patterns to guide the generation of inputs more likely to trigger vulnerabilities, without undermining the role of coverage feedback during the initial exploration phase. Existing solutions prove inadequate here as well. One approach, similar to KernelGPT, uses large language models to directly generate test inputs. However, this approach suffers from low generation efficiency and difficulty ensuring input validity. Another approach, exemplified by Snowplow, trains neural networks on collected data to guide mutation. Unfortunately, thousands of crash samples constitute a relatively small training set for neural networks, and the highly structured nature of these program fragments combined with substantial parameter noise renders them poorly suited for direct neural network training.

Our approach addresses these challenges through a three-phase methodology. First, we collect crash samples from tools like Syzbot and preprocess them using large language models to extract generalized vulnerability patterns in the form of Syzlang templates. Second, we design a Trie-based data structure to represent these patterns. Since vulnerability triggering often involves combinations of multiple system calls and their parameters, we model this as a high-dimensional Markov decision process. This formulation allows us to perform random walks over the Trie structure to generate novel test inputs that inherit the structural properties of past crashes. Third, we combine coverage feedback with crash-derived guidance. During the early phase of fuzzing, coverage feedback dominates the guidance mechanism. As coverage growth slows, we gradually increase the influence of crash-derived patterns, thereby sustaining effective guidance and enabling continued vulnerability discovery throughout extended fuzzing campaigns.

To evaluate our approach, we conduct experiments on the Linux 5.15 kernel. We extract the Trie structure from crash data associated with this kernel version and compare our method against state-of-the-art fuzzers including Syzkaller, KernelGPT, and Snowplow. Specifically, we first run baseline fuzzers for one month, then enable crash-guided fuzzing for an additional five days, comparing the number of discovered vulnerabilities and coverage growth with and without crash-derived guidance. The experimental results demonstrate that our method significantly improves vulnerability discovery efficiency during prolonged campaigns. In the 30+5 day experiments, our approach discovers 3.3 times more vulnerabilities and achieves 3.2 times more coverage growth during the final five days compared to the state-of-the-art baseline. Moreover, our method uncovers several deep vulnerabilities that would require hundreds of days or longer to discover using coverage-only approaches.

In summary, this paper makes the following contributions. First, we propose the first approach to leverage historical crash samples to guide kernel fuzzing, thereby overcoming the limitations of relying solely on coverage feedback. Second, we design a Trie-based representation for vulnerability patterns that effectively extracts and utilizes the information contained in crash samples to guide test input generation. Third, we conduct extensive experiments on Linux kernel fuzzing, demonstrating the effectiveness of our approach and showing significant improvements in vulnerability discovery efficiency, particularly during long-term fuzzing campaigns.